\begin{helper}
Fix natural numbers \(p,m,t\) and a prefix tuple
\[
  pref:\{0,\ldots,m-1\}\to \mathbb F_2^p\times\mathbb F_2^p.
\]
The number of possible last pairs \((a,b)\in\mathbb F_2^p\times\mathbb F_2^p\)
such that appending \(a\) to the prefix first-coordinate list gives rank \(t\)
and
\[
  \langle a_{\ell+1},b\rangle=1
  \quad\text{for every zero-based }\ell<m+1
\]
is at most
\[
  2^p\,2^{p-t}.
\]
As in \(\noderef{F2BadTupleLastBChoicesBound}\), Lean computes the rank after
appending \(a\) by appending the dummy pair \((a,0)\), since the rank depends
only on first coordinates.
\end{helper}
\begin{proof}
Partition the counted last pairs according to their first coordinate \(a\).
There are exactly \(2^p\) possible choices for \(a\in\mathbb F_2^p\). For a
fixed \(a\), \(\noderef{F2BadTupleLastBChoicesBound}\) bounds the number of
possible second coordinates \(b\) satisfying the displayed dot-product
equations and the rank-\(t\) condition by \(2^{p-t}\). Therefore the total
number of counted last pairs is at most the sum, over all \(2^p\) choices of
\(a\), of \(2^{p-t}\). This sum is
\[
  2^p\,2^{p-t}.
\]
\end{proof}
